% ============================================================
%  tikz_fig_adas_lr.tex
%  長期の総需要・総供給（AD-AS）分析
%
%  【単体コンパイル】
%    uplatex  tikz_fig_adas_lr.tex && dvipdfmx tikz_fig_adas_lr.dvi
%
%  【Beamer への組み込み（my_preamble.tex に standalone 追加済みの場合）】
%    \input{tikz/tikz_fig_adas_lr}
%
%  曲線の仕様:
%    AD  : P = 10/Y + 1.5   （右下がり双曲線）
%    LRAS: Y = Y*  = 5.0    （垂直線・潜在 GDP）
%    均衡: (Y*, P*) = (5, 3.5)
%
%  コメントアウトで以下も用意:
%    SRAS : 短期総供給曲線（右上がり）
%    AD'  : 総需要曲線のシフト後（ADシフトの例）
% ============================================================
\documentclass[tikz, dvipdfmx]{standalone}
\usepackage{tikz}
\usetikzlibrary{arrows.meta}
\definecolor{gdpblue} {RGB}{30,  90,  180}
\definecolor{trendred}{RGB}{190, 40,  40}
\definecolor{srasgreen}{RGB}{30, 140, 60}

\begin{document}
\begin{tikzpicture}[
  >=Stealth,
  thick,
  % 1 unit = 1 cm; 図全体の描画域はおよそ 9.5 cm × 7.5 cm
]

% ── 軸 ──────────────────────────────────────────────────────
\draw[->] (0,0) -- (9.8,0)
  node[right, font=\small] {$Y$ (Output)};
\draw[->] (0,0) -- (0,7.8)
  node[above, font=\small] {$P$ (Price level)};

% ── 長期総供給曲線 LRAS（垂直線） ────────────────────────
%    Y = Y*（潜在 GDP）
\draw[trendred, very thick] (5,0) -- (5,7.2)
  node[above, font=\small, trendred] {LRAS};

% ── 総需要曲線 AD（右下がり） ─────────────────────────────
%    P = 10/Y + 1.5
%    Y*=5 のとき P* = 10/5 + 1.5 = 3.5
%    domain: P=7.4 → Y≈1.71,  P→1.5（漸近）→ Y→∞
\draw[gdpblue, very thick]
  plot[domain=1.72:9.3, samples=120, smooth] (\x, {10/\x + 1.5})
  node[right, font=\small, gdpblue] {AD};

% ── 均衡点 (Y*, P*) = (5, 3.5) ────────────────────────────
\fill[black] (5,3.5) circle[radius=2.5pt];

% ── P* への水平点線 ──────────────────────────────────────
\draw[dashed, gray!70, thin]
  (0,3.5) node[left, font=\small, black] {$P^*$} -- (5,3.5);

% ── Y* のラベル（横軸下）────────────────────────────────
\node[below, font=\small] at (5,0) {$Y^*$};


% ════════════════════════════════════════════════════════════
%  以下はコメントアウト:  必要に応じて有効化
% ════════════════════════════════════════════════════════════

%% ── (A) 短期総供給曲線 SRAS（右上がり）──────────────────
%%    短期均衡を示したい場合に有効化
%% \draw[srasgreen, very thick]
%%   plot[domain=0.3:9, samples=2, smooth] (\x, {0.7*\x + 0.0})
%%   node[right, font=\small, srasgreen] {SRAS};

%% ── (B) AD シフト（例: 拡張的財政・金融政策） ────────────
%%    AD が右上にシフト → 長期では P のみ上昇, Y は不変
%%
%%    AD' : P = 14/Y + 1.5   （AD を上方シフト）
%%    新均衡: (Y*, P**) = (5, 14/5+1.5) = (5, 4.3)
%%
%% % AD' 曲線
%% \draw[gdpblue!50, very thick, dashed]
%%   plot[domain=1.72:9.3, samples=120, smooth] (\x, {14/\x + 1.5})
%%   node[right, font=\small, gdpblue!70] {AD$'$};
%%
%% % 新均衡点
%% \fill[black] (5,4.3) circle[radius=2.5pt];
%%
%% % P** への水平点線
%% \draw[dashed, gray!70, thin]
%%   (0,4.3) node[left, font=\small, black] {$P^{**}$} -- (5,4.3);
%%
%% % AD から AD' へのシフト矢印
%% \draw[->, thick, gdpblue!60]
%%   (3.5, {10/3.5+1.5}+0.15) -- (3.5, {14/3.5+1.5}-0.15);

\end{tikzpicture}
\end{document}
